\section{Pengenalan Composer dan Dependency Management}

\textbf{Composer} adalah package manager untuk PHP --- mirip npm untuk JavaScript. Composer mengelola library pihak ketiga (\textit{dependency}) yang digunakan proyek Anda: instalasi, update, dan autoloading. File \code{composer.json} mendefinisikan dependency proyek; file \code{composer.lock} mengunci versi tepat yang diinstal. Direktori \code{vendor/} berisi library yang terunduh. Perintah utama: \code{composer init} (membuat proyek), \code{composer require vendor/package} (instal library), \code{composer install} (instal dari lock file), \code{composer update} (update dependency) \cite{phpRightWay}.

\begin{webcode}[language=PHP, caption={Contoh composer.json dan penggunaan library}]
// composer.json
{
  "name": "myapp/toko",
  "require": {
    "php": ">=8.0",
    "vlucas/phpdotenv": "^5.5",
    "phpmailer/phpmailer": "^6.8"
  },
  "autoload": {
    "psr-4": {
      "App\\": "src/"
    }
  }
}

// Terminal: composer install

// PHP: gunakan autoloader dan library
// require 'vendor/autoload.php';
// $dotenv = Dotenv\Dotenv::createImmutable(__DIR__);
// $dotenv->load();
// $db_host = $_ENV['DB_HOST'];
\end{webcode}

